\section{Experimental Setup}
\label{sec:setup}
\begin{table}[t]
\caption{Benchmarks. ``Contexts repeat'' means the same context appears in the training data with several outcomes, which removes the possibility of memorizing the realized outcome. Details in Appendix~\ref{app:bench}.}
\label{tab:bench}
\begin{center}
\footnotesize
\setlength{\tabcolsep}{4pt}
\resizebox{\linewidth}{!}{%
\begin{tabular}{lllllc}
\toprule
Name & Source & Stochasticity & Repeat & Actions & Model \\
\midrule
A1 Bouncing Shapes & FlowWM generator & normal bounce vs.\ reversal per collision & yes & none & 8M (joint 27--29M) \\
A2-div / A2-fix bounce & ours & $K \in \{2,4,8,16\}$ directions at a wall; asymmetric $K{=}2$ & no / partly & none & 8M \\
Push (branch risk) & ours & $\pm20^\circ$ (2 modes), ice $\{-60^\circ,0,60^\circ\}$ (3), absorbing pit & no & 2-D & 8M \\
B stochastic Push-T & DINO-WM data, re-simulated & hidden center of gravity, 4 modes & no & 10-D & 21M \\
BAIR robot pushing & \citet{ebert2017bair} & continuous, high-dimensional & no & 4-D & 21M \\
\bottomrule
\end{tabular}}
\end{center}
\end{table}
\textbf{Benchmarks} (Table~\ref{tab:bench}; Appendix~\ref{app:bench}). \emph{A1}: the public FlowWM generator, used unchanged \citep{porcher2026flowwm}: a square and a circle that, on hitting a wall, bounce normally or reverse their velocity with probability one half; $16{,}384$ training clips (one eighth of FlowWM's) from $5{,}324$ initial conditions with $3.08$ futures each on average, so contexts repeat; FlowWM's DINOv3 features, protocols and heatmap readout. \emph{A2}: a ball moving in a $64\times64$ frame whose direction after hitting a wall is drawn uniformly from $K$ equiangular directions (about $28\%$ of transitions are stochastic; no actions). A2-div redraws the background color, ball color, radius, speed and two distractor rectangles for every episode, so the nearest training context of a test context is $1.9\times$ farther away than the typical nearest neighbor within the training set and every training context carries an identifying key; A2-fix keeps the appearance fixed ($0.35\times$). Each setting has $40$k training frames and 400 held-out wall contexts; an asymmetric $K{=}2$ variant has mode probabilities $(0.8, 0.2)$. \emph{Push}: the same renderer with a 2-D push action deflected by $\pm20^\circ$ in the open region (two modes) and by $\{-60^\circ, 0, 60^\circ\}$ on an ice patch beside an absorbing pit (three modes); in the \emph{branch-risk task} the path to the goal runs along the pit, so one of the three ice modes leads into it. \emph{B}: the Push-T action sequences of DINO-WM \citep{zhou2025dinowm} replayed in the simulator with the T block's center of gravity hidden at one of four positions per episode (mass and friction have no effect under quasi-static pushing); $2{,}400$ training episodes; a multi-future test set of 128 pre-contact starting states $\times$ 64 futures $\times$ 16 steps. \emph{BAIR}: $64\times64$ frames with 4-D actions; $4{,}096$ training and 256 evaluation sequences with 16 samples each.

\textbf{Metrics.} \emph{Coverage} and \emph{precision}: 16 samples per held-out stochastic context are decoded by a position readout (test error $0.13$\,px; Appendix~\ref{app:readout}) and assigned to the nearest enumerated mode within a tolerance; coverage is the fraction of modes hit, and precision is the fraction of samples that land on a mode (an ideal 16-sample sampler reaches coverage $1 - (1-1/K)^{16}$). \emph{Memorization index}: the same share, on the realized mode of 200 training contexts (chance $1/K$; \emph{excess memorization} is $(\text{index} - 1/K)/(1-1/K)$). \emph{kNN precision/recall} \citep{kynkaanniemi2019improved} on PCA-reduced tokens; \emph{deterministic-transition error}: the readout error on non-stochastic transitions. \emph{FlowWM precision/recall error} (A1, B): with twice as many rollouts as true futures per starting state, the mean distance from each rollout to the nearest true future (precision), from each true future to the nearest rollout (recall), and their harmonic mean (F1); L1 over the shape centers on A1, L2 over the eight T vertices on B. \emph{Manifold distance}: the distance from a predicted token map to the nearest real one, divided by the median nearest-neighbor distance among real maps ($1$ = on the data). \emph{Best-of-$n$}: the readout (or token) error of the sample nearest the truth; \emph{spread}: the mean pairwise distance among the samples. On BAIR we also report token RMS errors, the spread-to-error ratio, and the Fr\'echet distance and kNN recall of pooled tokens. \emph{Latency} is measured per predicted step at batch size 1; \emph{planning} reports success rate, pit rate and wall-clock time per MPC step.

\textbf{Protocol.} All models use AdamW (learning rate $2\times10^{-4}$, $\beta_2 = 0.9$), batch size 128, an exponential moving average of the weights, and 12k training steps unless stated otherwise (8k for the budget grid, the asymmetric variant and per-step matching; 36k for iMF; 24k for Shortcut); parameter counts match within $\pm4\%$. Main comparisons use three seeds, grids and follow-up runs two, and ablations one; uncertainty is quantified with bootstrap intervals over contexts and, for Push-T, a hierarchical bootstrap over seeds and starting states (Appendix~\ref{app:stats}). All experiments ran on two RTX A5000 GPUs, about 300 GPU-hours in total. The predictions for the $K{=}2,16$ grid, the energy-score follow-up runs and the fed-back-state ablation were written down before those runs and are reproduced in the appendix, including the ones that missed.
